\section{Background}

%This section describes the definition of our task, and a baseline approach with state-of-the-art pre-trained models for multi-class classification.

\subsection{Task: Few-Shot Intent Detection}
\label{subsec:task-def}

Given a user utterance $u$ at every turn in a goal-oriented dialog system, an intent detection model $I(u)$ aims at predicting the speaker's intent:
\begin{equation}
    I(u) = c,
\end{equation}
where $c$ is one of pre-defined $N$ intent classes $\mathbf{C}=\{C_1, C_2, \ldots, C_N\}$, or is categorized as OOS.
The OOS category corresponds to user utterances whose requests are not covered by the system.
In other words, any utterance can be OOS as long as it does {\it not} fall into any of the $N$ intent classes, so the definition of OOS is different depending on $\mathbf{C}$.

\paragraph{Balanced $K$-shot learning}
In a few-shot learning scenario, we have a limited number of training examples for each class, and we assume that we have $K$ examples for each of the $N$ classes in our training data.
In other words, we have $N\cdot K$ training examples in total.
We denote the $i$-th training example from the $j$-th class $C_j$ as $e_{j,i}\in E$, where $E$ is the set of the examples.
$K$ is typically 5 or 10.


\subsection{Multi-Class Classification}
\label{subsec:classifier}

The goal is to achieve high accuracy both for the intent classification and OOS detection.
One common approach to this task is using a multi-class classification model.
% More specifically, we can leverage a general pre-trained model as a basic building block, since it is reasonable to start from a pre-trained model especially for few-shot learning.
Specifically, to get a strong baseline for the few-shot learning use case, one can leverage a pre-trained model as transfer learning, which has been shown to achieve state-of-the-art results on numerous natural language processing tasks.
We use BERT~\citep{bert,roberta} as a text encoder:
\begin{equation}
    \label{eq:bert}
    h = \mathrm{BERT}(u) \in\mathbb{R}^{d},
\end{equation}
where $h$ is a $d$-dimensional output vector corresponding to the special token {\tt[CLS]} as in the follwing input format: [{\tt[CLS]}, $u$, {\tt[SEP]}].\footnote{The format of these special tokens is different in RoBERTa, but we use the original BERT's notations.}

To handle the intent classification and the OOS detection, we apply the threshold-based strategy in \citet{oos-intent}, to the softmax output of the $N$-class classification model~\citep{softmax_conf}:
\begin{equation}
    \label{eq:softmax}
    p(c | u) = \mathrm{softmax}(W h + b) \in\mathbb{R}^{N},
\end{equation}
where $W\in\mathbb{R}^{N\times d}$ and $b\in\mathbb{R}^{N}$ are the classifier's model parameters.
The classification model is trained by cross-entropy loss with the ground-truth intent labels of the training examples.
At inference time, we first take the class $C_j$ with the largest value of $p(c=C_j|u)$, then output $I(u)=C_j$ if $p(c=C_j|u) > T$, where $T\in[0.0, 1.0]$ is a threshold to be tuned, and otherwise we output $I(u)=\mathrm{OOS}$.


\begin{table*}[t]

 \begin{center}
{\small
 \resizebox{1.0\linewidth}{!}{
    \begin{tabular}{l|l|l|l}
    
    & Input utterance & Utterance to be compared & Label \\ \hline
    (a) & i need you to send 500 dollars from & help me move my money please & pos. \\
        & my high tier account to my regular checking & & \\ \hdashline
    (b) & i would like to know when the bill is due & i need to know the amounts due for my utilities and cable bills & neg. \\ \hline
    %(c) & what is saas marketing & why would i have been locked out of my own account & neg. \\ \hline

    (c) & And can you tell me any of the names and & Are you able to inform me of any name or address? & pos. \\
        & addresses? Annie considered. & & \\ \hdashline
    (d) & It's Sunday, what channel is this? & It's Sunday, can you change the channel? & neg. \\ \hdashline
    (e) & I want to go back to Marguerite. & I never want to return to Marguerite. & neg. \\ \hline

    \end{tabular}}
%}
    \caption{Training examples for our model. The first two examples ((a)--(b)) come from the CLINC150 dataset~\citep{oos-intent}, and the other three examples ((c)--(e)) come from the MNLI dataset~\citep{mnli}.}
    \label{tb:examples}
    \vspace{-1em}
}
  \end{center}
\end{table*}

\subsection{Nearest Neighbor Classification}

As the fundamental building block of our proposed method, we also review nearest neighbor classification (i.e., $k$-nearest neighbors ($k$NN) classification with $k=1$), a simple and well-established concept for classification~\citep{knn-nips,knn-classifier}.
The basic idea is to classify an input into the same class of the most relevant training example based on a certain metric.

In our task, we formulate a nearest neighbor classification model as the following:
\begin{equation}
    \label{eq:example-based}
    I(u) = \mathrm{class}\left(\argmax{e_{j, i}\in E} S(u, e_{j, i}) \right),
\end{equation}
where $\mathrm{class}(e_{j, i})$ is a function returning the intent label (class) of the training example $e_{j, i}$, and $S$ is a function that estimates some relevance score between $u$ and $e_{j, i}$.
To detect OOS, we can also use the uncertainty-based strategy in Section~\ref{subsec:classifier}; that is, we take the output label from Equation~(\ref{eq:example-based}) if the corresponding relevance score is greater than a threshold, and otherwise we output $I(u)=\mathrm{OOS}$.



\section{Proposed Method}
\label{subsec:method}


This section first describes how to directly model inter-utterance relations in our nearest neighbor classification scenario.
We then introduce a binary classification strategy by synthesizing pairwise examples, and propose a seamless transfer of NLI.
Finally, we describe how to speedup our method's inference process.

\subsection{Deep Pairwise Matching Function}
The objective of $S(u, e_{j, i})$ in Equation~(\ref{eq:example-based}) is to find the best matched utterance from the training set $E$, given the input utterance $u$.
The typical methodology is to embed each data example into a vector space and (1) use an off-the-shelf distance metric to perform a similarity search~\citep{knn-classifier} or (2) learn a distant metric between the embeddings~\citep{relation_net}.
However, as shown in Figure~\ref{fig:tsne-main}, the text embedding methods do not discriminate the OOS examples well enough.

To model fine-grained relations of utterance pairs to distinguish in-domain and OOS intents, we propose to formulate $S(u, e_{j, i})$ as follows:
\begin{eqnarray}
    h &=& \mathrm{BERT}(u, e_{j, i}) \in\mathbb{R}^{d}, \label{eq:NLI} \\
    S(u, e_{j, i}) &=& \sigma(W \cdot h + b) \in\mathbb{R}, \label{eq:sigmoid}
\end{eqnarray}
where $\mathrm{BERT}$ is the same model in Equation~(\ref{eq:bert}), except that we follow a different input format to accommodate pairs of utterances: [{\tt[CLS]}, $u$, {\tt[SEP]}, $e_{j, i}$, {\tt[SEP]}].
$\sigma$ is the sigmoid function, and $W\in\mathbb{R}^{1\times d}$ and $b\in\mathbb{R}$ are the model parameters.
We can interpret our method as wrapping both the embedding and matching functions into the paired encoding with the deep self-attention in BERT (Equation~(\ref{eq:NLI})) along with the discriminative model (Equation~(\ref{eq:sigmoid})).
It has been shown that the paired text encoding is crucial in capturing complex relations between queries and documents in document retrieval~\citep{watanabe-query,sem-bert,graph-retriever}.


\subsection{Discriminative Training}
\label{subsec:pos_neg}

We train the matching model $S(u, e_{j, i})$ as a binary classifier, such that $S(u, e_{j, i})$ is closed to 1.0 if $u$ belongs to the same class of $e_{j, i}$, and otherwise closed to 0.0.
The model is trained by a binary cross-entropy loss function.

\paragraph{Positive examples}
To create positive examples, we consider all the possible ordered pairs within the same intent class: $(e_{j, i}, e_{j, \ell})$ such that $i \neq \ell$.
We thus have $N\times K\times(K-1)$ positive examples in total.
Table~\ref{tb:examples} (a) shows a positive example created from an intent, ``transfer,'' in a banking domain.

\paragraph{Negative examples}
For negative examples, we consider all the possible ordered pairs across any two different intent classes: $(e_{j, i}, e_{o, \ell})$ such that $j \neq o$.
We thus have $K^2\times N\times(N-1)$ negative examples in total, and this number is in general greater than that of the positive examples.
%(i.e. $\mathcal{O}(K^2N^2) > \mathcal{O}(K^2N)$).
Table~\ref{tb:examples} (b) shows a negative example, where the input utterance comes from the intent, ``bill due'', and the paired sentence from another intent, ``bill balance''. 


\subsection{Seamless Transfer from NLI}
\label{subsec:nli_transfer}

A key characteristic of our method is that we seek to model the relations between the utterance pairs, instead of explicitly modeling the intent classes.
To mitigate the data scarcity setting in few-shot learning, we consider transferring another inter-sentence-relation task.

This work focuses on NLI; the task is to identify whether a hypothesis sentence can be entailed by a premise sentence~\citep{deep-nli}.
We treat the NLI task as a binary classification task: {\tt entailment} (positive) or {\tt non-entailment} (negative).\footnote{A widely-used format is a three-way classification task with {\tt entailment}, {\tt neutral}, and {\tt contradiction}, but we merge the latter two classes into a single {\tt non-entailment} class.}
We first pre-train our model with the NLI task, where the premise sentence corresponds to the $u$-position, and the hypothesis sentence corresponds to the $e_{j, i}$-position in Equation~(\ref{eq:NLI}).
Note that it is not necessary to modify the model architecture since the task format is consistent, and we can train the NLI model solely based on existing NLI datasets.
Once the NLI model pre-training is completed, we fine-tune the NLI model with the intent classification training examples described in Section~\ref{subsec:pos_neg}.
This allows us to transfer the NLI model to any intent detection datasets seamlessly.

\paragraph{Why NLI?}
The NLI task has been actively studied, especially since the emergence of large scale datasets~\citep{snli,mnli}, and we can directly leverage the progress.
%The existing datasets are designed to describe a single event in a sentence~\citep{deep-nli}, which is well-aligned with the intent detection task.
Moreover, recent work is investigating cross-lingual NLI~\citep{eriguchi-crosslingual,xnli}, and this is encouraging to consider multilinguality in future work.
On the other hand, while we can find examples relevant to the intent detection task, as shown in Table~\ref{tb:examples} ((c), (d), and (e)), we still need the few-shot fine-tuning.
This is because a domain mismatch still exists in general, and perhaps more importantly, our intent detection approach is not exactly modeling NLI.
%Additional notes are available in Appendix~A.

\paragraph{Why not other tasks?}
There are other tasks modeling relationships between sentences.
Paraphrase~\citep{paranmt} and semantic relatedness~\citep{semeval} tasks are such examples.
It is possible to automatically create large-scale paraphrase datasets by machine translation~\citep{ppdb}.
However, our task is not a paraphrasing task, and creating negative examples is crucial and non-trivial~\citep{selectional-preference}.
In contrast, as described above, the NLI setting comes with negative examples by nature.
The semantic relatedness (or textual similarity) task is considered as a coarse-grained task compared to NLI, as discussed in the previous work~\citep{jmt}, in that the task measures semantic or topical relatedness.
This is not ideal for the intent detection task, because we need to discriminate between topically similar utterances of different intents.
In summary, the NLI task well matches our objective, with access to the large datasets.



\subsection{A Joint Approach with Fast Retrieval}
\label{subsec:joint}

The number of model parameters of the multi-class classification model in Section~\ref{subsec:classifier} and our model in Section~\ref{subsec:method} is almost the same when we use the same pre-trained models.
%\footnote{The actual difference is $(N-1)(d+1)$, but this is negligible compared with the model size of BERT.}
% However, our example-based method has an inference-time bottleneck in Equation~(\ref{eq:NLI}), where we need the paired BERT encoding with all the training examples, given an input utterance $u$.
However, our example-based method has an inference-time bottleneck in Equation~(\ref{eq:NLI}), where we need to compute the BERT encoding for all $N\times K$ $(u, e_{j,i})$ pairs. 

We follow common practice in document retrieval to reduce the inference-time bottleneck~\citep{sem-bert,graph-retriever}, by introducing a fast text retrieval model to select a set of top-$k$ examples $E_k$ from the training set $E$, based on its retrieval scores.
We then replace $E$ in Equation~(\ref{eq:example-based}) with the shrunk set $E_k$.
The cost of the paired BERT encoding is now constant, regardless the size of $E$.
Either TF-IDF~\citep{drqa} or embedding-based retrieval~\citep{faiss,phrase-embedding-index,latent-retrieval} can be used for the first step.
We use the following fast $k$NN.

\paragraph{Faster $k$NN}
%Our target text is much shorter than the documents used in open-domain question answering.
%Therefore it could perform well to separately encode query and target embeddings, without the deep self attention in Equation~(\ref{eq:NLI}).
As a baseline and a way to instantiate our joint approach, we use Sentence-BERT (SBERT)~\citep{sbert} to separately encode $u$ and $e_{j, i}$ ($x\in\{u, e_{j, i}\}$) as follows:
\begin{eqnarray}
    \label{eq:sbert}
    v(x) &=& \mathrm{SBERT}(x) \in\mathbb{R}^{d},
\end{eqnarray}
where the input text format is identical to that of BERT in Equation~(\ref{eq:bert}).
SBERT is a BERT-based text embedding model, fine-tuned by siamese networks with NLI datasets.
Thus both our method and SBERT transfer the NLI task in different ways.

% We then replace $S(u, e_{j, i})$ in Equation~(\ref{eq:sigmoid}) with the cosine similarity between $v(u)$ and $v(e_{j, i})$.
Cosine similarity between $v(u)$ and $v(e_{j, i})$ then replaces $S(u, e_{j, i})$ in Equation~(\ref{eq:sigmoid}).
To get a fair comparison, instead of using the encoding vectors produced by the original SBERT, we fine-tune SBERT with our intent training examples described in Section~\ref{subsec:pos_neg}.
The cosine similarity is symmetric, so we have half the training examples.
We use the pairwise cosine-based loss function in \citet{sbert}.
After the model training, we pre-compute $v(e_{j, i})$ for fast retrieval.

%\subsubsection{Prototypical Networks}
%Describe this later!


%\paragraph{Hierarchical Attention Prototypical Networks}

